\begin{theorem}[Eisenstein series for $\SL_2(\mathbb{Z})$ is a modular form of weight $k$]
\label{theorem:eisenstein_series_for_SL_2_Z_is_a_modular_form_with_fourier_expansion}
\TODO{AI generated, verify}
Let $k \geq 3$ be an integer. The \CrefAndHyperrefIfExist{definition:eisenstein_series_of_weight_k_for_a_congruence_subgroup_of_SL_2_Z}{unnormalized Eisenstein series $G_k(z)$} for $\SL_2(\mathbb{Z})$ defines a \CrefAndHyperrefIfExist{definition:modular_form_and_cusp_form_of_weight_k_with_respect_to_a_congruence_subgroup_of_SL_2_Z}{modular form} of weight $k$ for $\SL_2(\mathbb{Z})$. Its \CrefAndHyperrefIfExist{definition:fourier_expansion_of_a_periodic_holomorphic_function_on_the_upper_half_plane}{Fourier expansion} at the \CrefAndHyperrefIfExist{definition:cusps_of_a_congruence_subgroup_of_SL_2_Z}{cusp} $\infty$ is given by
$$G_k(z) = 2\zeta(k) + 2\frac{(-2\pi i)^k}{(k-1)!} \sum_{n=1}^{\infty} \sigma_{k-1}(n) \, e^{2\pi i n z},$$
where $\zeta(s)$ is the \CrefAndHyperrefIfExist{definition:riemann_zeta_function}{Riemann zeta function} and $\sigma_{k-1}(n) = \sum_{d|n} d^{k-1}$ is the \hldef{sum of $(k-1)$-st powers of divisors function} \TODO{AI generated, reference needed: definition of divisor sum function $\sigma_k(n) = \sum_{d|n} d^k$}.

Consequently, the \CrefAndHyperrefIfExist{definition:eisenstein_series_of_weight_k_for_a_congruence_subgroup_of_SL_2_Z}{normalized Eisenstein series $E_k(z)$} has the Fourier expansion
$$E_k(z) = 1 + \frac{(-2\pi i)^k}{(k-1)! \, \zeta(k)} \sum_{n=1}^{\infty} \sigma_{k-1}(n) \, e^{2\pi i n z}.$$
\end{theorem}

\begin{proof}
\TODO{AI generated, verify}
We verify each of the defining properties of a modular form.

\begin{enumerate}
  \item We show that $G_k(z)$ \CrefAndHyperrefIfExist{definition:absolute_convergence_of_a_series_in_module_over_a_commutative_ring_with_an_absolute_value_valued_in_an_ordered_semiring}{converges absolutely} and \CrefAndHyperrefIfExist{definition:uniform_convergence_of_a_sequence_of_functions_from_a_set_to_an_extended_metric_space_valued_in_ordered_commutative_monoid_with_adjoined_infinity}{uniformly} on compact subsets of $\mathbb{H}$. Let $K \subset \mathbb{H}$ be compact. For each $z = x + iy \in \mathbb{H}$,
  $$|mz + n|^2 = |(mx+n) + imy|^2 = (mx + n)^2 + (my)^2 = m^2(x^2 + y^2) + 2mnx + n^2.$$
  This is a positive definite quadratic form in $(m,n)$ with matrix
  $$\begin{pmatrix} |z|^2 & x \\ x & 1 \end{pmatrix}.$$
  Its minimum eigenvalue
  $$\lambda_{\min}(z) = \frac{|z|^2 + 1 - \sqrt{(|z|^2 - 1)^2 + 4x^2}}{2}$$
  is positive because $y = \Im(z) > 0$ (the discriminant satisfies $(|z|^2-1)^2 + 4x^2 = (|z|^2+1)^2 - 4y^2 < (|z|^2+1)^2$, so $\lambda_{\min}(z) > 0$). The function $\lambda_{\min} : \mathbb{H} \to \mathbb{R}_{>0}$ is continuous, and $K$ is compact, so
  $$c := \min_{z \in K} \lambda_{\min}(z) > 0.$$
  By the \CrefAndHyperrefIfExist{theorem:spectral_theorem_for_normal_operators_on_a_finite_dimensional_inner_product_space}{spectral theorem}, for every $z \in K$ and every $(m,n) \in \mathbb{R}^2$,
  $$|mz + n|^2 \geq \lambda_{\min}(z) \cdot (m^2 + n^2) \geq c\,(m^2 + n^2).$$
  Consequently,
  $$\frac{1}{|mz + n|^k} \leq \frac{1}{c^{k/2}} \cdot \frac{1}{(m^2 + n^2)^{k/2}}.$$
  By \Cref{definition:series_indexed_by_a_countable_set}, the double series $\sum_{(m,n) \neq (0,0)} (m^2 + n^2)^{-k/2}$ is defined by the supremum of its finite partial sums. The function $(x,y) \mapsto (x^2+y^2)^{-k/2}$ is nonincreasing in each variable separately on $[0,\infty)^2$. By \Cref{theorem:comparison_of_double_series_to_double_integrals_of_nonnegative_functions} with $m_0 = n_0 = 1$, we have
	  $$\sum_{m=2}^{\infty} \sum_{n=2}^{\infty} (m^2 + n^2)^{-k/2} \leq \iint_{[1,\infty)^2} (x^2 + y^2)^{-k/2}\,dx\,dy.$$
	  \CrefIfExists{definition:lebesgue_integral_of_measurable_functions_into_the_real_line_extended_by_infinities}
	  The remaining terms of the double series (those with $m=0$, $m=1$, $n=0$, or $n=1$) are bounded by the convergent single series $\sum_{m=1}^\infty m^{-k} = \zeta(k) < \infty$, and similarly for $n$.

	  To evaluate the double integral, apply \Cref{theorem:change_of_variables_for_lebesgue_integrals_on_Rn} to the \CrefAndHyperrefIfExist{definition:polar_coordinates_on_R2}{polar coordinate map} $\Phi(r,\theta) = (r\cos\theta, r\sin\theta)$. On the region $U = \{(r,\theta) \in (0,\infty) \times (0,\pi/2) : r\cos\theta \ge 1,\; r\sin\theta \ge 1\}$, the map $\Phi$ is a $C^1$ diffeomorphism onto $[1,\infty)^2$, and $|\det D\Phi(r,\theta)| = r$. Hence
	  \begin{align*}
	  \iint_{[1,\infty)^2} (x^2+y^2)^{-k/2}\,dx\,dy
	  &= \iint_U r^{-k} \cdot r\,dr\,d\theta \\
	  &\le \int_0^{\pi/2} \int_1^\infty r^{1-k}\,dr\,d\theta
	  = \frac{\pi}{2} \int_1^\infty r^{1-k}\,dr,
	  \end{align*}
	  which is finite for $k > 2$ because $\int_1^\infty r^{1-k}\,dr = 1/(k-2)$. Thus the double series converges. By the \CrefAndHyperrefIfExist{theorem:weierstrass_m_test_for_uniform_convergence_of_series_of_functions_from_sets_to_banach_spaces_over_a_absolute_valued_ring_valued_in_the_reals}{Weierstrass $M$-test}, the series defining $G_k(z)$ \CrefAndHyperrefIfExist{definition:absolute_convergence_of_a_series_in_module_over_a_commutative_ring_with_an_absolute_value_valued_in_an_ordered_semiring}{converges absolutely} and \CrefAndHyperrefIfExist{definition:uniform_convergence_of_a_sequence_of_functions_from_a_set_to_an_extended_metric_space_valued_in_ordered_commutative_monoid_with_adjoined_infinity}{uniformly} on $K$.

  \item We show thta $G_k(z)$ is Holomorphic on $\mathbb{H}$. Each term $(mz+n)^{-k}$ is a holomorphic function of $z$ on $\mathbb{H}$ (it is a rational function with a simple pole at $z = -n/m$ for $m \neq 0$, which lies on the real axis, not in $\mathbb{H}$). The series converges uniformly on compact subsets of $\mathbb{H}$ by the previous step. By \Cref{theorem:the_weierstrass_convergence_theorem_stating_that_the_uniform_limit_on_compact_subsets_of_a_sequence_of_holomorphic_functions_is_itself_holomorphic_with_all_derivatives_converging_uniformly_on_compact_subsets}, the limit function $G_k(z)$ is holomorphic on $\mathbb{H}$.

  \item We verify that $G_k$ is \CrefAndHyperrefIfExist{definition:weakly_modular_form_with_respect_to_a_subgroup_of_SL_2_Z}{weakly modular}. For $\gamma = \begin{pmatrix} a & b \\ c & d \end{pmatrix} \in \SL_2(\mathbb{Z})$, we compute:
  \begin{align*}
    G_k(\gamma z) &= \sum_{(m,n) \neq (0,0)} \frac{1}{\bigl(m \cdot \frac{az+b}{cz+d} + n\bigr)^k} \\
    &= \sum_{(m,n) \neq (0,0)} \frac{(cz+d)^k}{\bigl(m(az+b) + n(cz+d)\bigr)^k} \\
    &= (cz+d)^k \sum_{(m,n) \neq (0,0)} \frac{1}{\bigl((ma+nc)z + (mb+nd)\bigr)^k}.
  \end{align*}
  Since $\gamma$ is an \CrefAndHyperrefIfExist{definition:invertible_matrix_over_a_ring}{invertible} $2 \times 2$ matrix of integers, the map $(m,n) \mapsto (m,n)\gamma = (ma+nc, mb+nd)$ is a permutation of $\mathbb{Z}^2 \setminus \{(0,0)\}$. Hence the sum over $(m,n)$ is identical to the sum over $(m',n') = (m,n)\gamma$, and we obtain
  $$G_k(\gamma z) = (cz+d)^k \sum_{(m',n') \neq (0,0)} \frac{1}{(m'z + n')^k} = (cz+d)^k G_k(z).$$
  Thus $G_k$ satisfies the modular transformation law of weight $k$.

  \item We show that $G_k(z)$ is \CrefAndHyperrefIfExist{definition:holomorphic_at_infinity_for_a_meromorphic_function_on_the_upper_half_plane}{holomorphic at $\infty$} and derive its Fourier series. By \Cref{lemma:weakly_modular_form_for_any_congruence_subgroup_of_SL_2_Z_is_periodic}, $G_k(z+1) = G_k(z)$ because $\begin{pmatrix} 1 & 1 \\ 0 & 1 \end{pmatrix} \in \SL_2(\mathbb{Z})$. Indeed,
  $$G_k(z+1) = \sum_{(m,n) \neq (0,0)} \frac{1}{(m(z+1) + n)^k} = \sum_{(m,n) \neq (0,0)} \frac{1}{(mz + (m+n))^k} = G_k(z)$$
  by the change of variables $(m,n) \mapsto (m, n-m)$.

  By \Cref{theorem:periodic_holomorphic_function_on_upper_half_plane_descends_to_punctured_disk_and_has_fourier_laurent_expansion}, $G_k$ has a Fourier expansion
  $$G_k(z) = \sum_{n=-\infty}^{\infty} a_n e^{2\pi i n z}.$$
  To compute the coefficients, we split the sum defining $G_k$ into the $m = 0$ and $m \neq 0$ parts.
  Absolute convergence of the double series (established above) implies, by \Cref{lemma:absolute_convergence_over_a_countable_index_set_is_unconditional}, that the sum is independent of the \CrefAndHyperrefIfExist{definition:injective_surjective_bijective_map_of_sets}{enumeration} of $\mathbb{Z}^2\setminus\{(0,0)\}$. Grouping the terms as $m=0$ then $m>0$ then $m<0$ specifies one such enumeration, so
  $$G_k(z) = \sum_{n \neq 0} \frac{1}{n^k} + \sum_{m=1}^{\infty} \sum_{n=-\infty}^{\infty} \frac{1}{(mz + n)^k} + \sum_{m=-\infty}^{-1} \sum_{n=-\infty}^{\infty} \frac{1}{(mz + n)^k}.$$
  The first sum is $2\sum_{n=1}^{\infty} n^{-k} = 2\zeta(k)$ (for $k$ even, this is well-known; for odd $k$ the sum still converges absolutely).

  For the second sum, we apply \Cref{theorem:lipschitz_summation_formula_for_reciprocal_powers_on_the_upper_half_plane}. For $m \geq 1$,
  $$\sum_{n=-\infty}^{\infty} \frac{1}{(mz + n)^k} = \frac{(-2\pi i)^k}{(k-1)!} \sum_{\ell=1}^{\infty} \ell^{k-1} e^{2\pi i \ell m z},$$
  and similarly for $m \leq -1$ with $z$ replaced by $-|m|z$. Summing over all $m \neq 0$ collects the coefficients into divisor sums:
  \begin{align*}
    G_k(z) &= 2\zeta(k) + \frac{(-2\pi i)^k}{(k-1)!} \sum_{m=1}^{\infty} \sum_{\ell=1}^{\infty} \ell^{k-1} e^{2\pi i \ell m z} + \frac{(-2\pi i)^k}{(k-1)!} \sum_{m=-\infty}^{-1} \sum_{\ell=1}^{\infty} \ell^{k-1} e^{2\pi i \ell m z} \\
    &= 2\zeta(k) + \frac{(-2\pi i)^k}{(k-1)!} \sum_{n=1}^{\infty} \left( \sum_{d|n} d^{k-1} \right) e^{2\pi i n z} \quad (\text{setting } n = \ell m) \\
    &= 2\zeta(k) + 2\frac{(-2\pi i)^k}{(k-1)!} \sum_{n=1}^{\infty} \sigma_{k-1}(n) e^{2\pi i n z}.
  \end{align*}
  (The factor of $2$ arises because the $m = -\ell$ terms contribute the same sum.) This is the claimed Fourier expansion. In particular, there are no negative powers of $e^{2\pi i z}$, so $G_k$ is holomorphic at $\infty$ by \Cref{definition:holomorphic_at_infinity_for_a_meromorphic_function_on_the_upper_half_plane}.
\end{enumerate}
  Having verified all four properties -- absolute convergence and uniform convergence on compacta yielding holomorphicity on $\mathbb{H}$, the modular transformation law, periodicity, and holomorphicity at $\infty$ with the explicit Fourier expansion -- we conclude that $G_k$ is a modular form of weight $k$ for $\SL_2(\mathbb{Z})$.
\end{proof}